\clearpage
\VerdierPhysicalPageAnchor{0070}%
\VerdierSemanticAnchor{ch1-2-2-3-continuation}%

\VerdierSemanticAnchor{ch1-2-2-3-eq3}
\begin{equation}
  \begin{aligned}
    \left(\int_{\varphi}^{\Pi}Y\right)^{\tau}
      &=\prod_{\varphi(\sigma)=\tau}Y^{\sigma},\\
    d_{\int_{\varphi}^{\Pi}Y}^{\tau,i}
      &\quad\text{由矩阵 }M(\tau,i,Y)\text{ 定义},\\[2pt]
    \left(\int_{\varphi}^{\Sigma}Y\right)^{\tau}
      &=\bigoplus_{\varphi(\sigma)=\tau}Y^{\sigma},\\
    d_{\int_{\varphi}^{\Sigma}Y}^{\tau,i}
      &\quad\text{由矩阵 }M(\tau,i,Y)\text{ 定义}.
  \end{aligned}
  \tag{2.2.3.3}
\end{equation}

\VerdierSemanticAnchor{ch1-2-2-3-pi-extension}
函子 $\int_{\varphi}^{\Pi}$ 称为函子 $\int_\varphi$ 的
$\Pi$-\emph{扩张}。

\VerdierSemanticAnchor{ch1-2-2-3-sigma-extension}
函子 $\int_{\varphi}^{\Sigma}$ 称为函子 $\int_\varphi$ 的
$\Sigma$-\emph{扩张}。下文将使用这两种扩张。

\VerdierSemanticAnchor{ch1-2-2-3-natural-map}
此外，还有函子的自然态射
$\int_{\varphi}^{\Sigma}\to\int_{\varphi}^{\Pi}$，它对应于从直和到乘积的
函子性态射。

\VerdierSemanticAnchor{ch1-2-2-4}
\noindent\textbf{2.2.4.}\quad
设 $[n]\xrightarrow{\varphi}[m]\xrightarrow{\psi}[r]$ 为两个映射，并为明确
起见假定 $C$ 中可数乘积可表示。于是有关系

\VerdierSemanticAnchor{ch1-2-2-4-eq1}
\begin{equation}
  \int_{\psi}^{\Pi}\int_{\varphi}^{\Pi}
  =\int_{\psi\varphi}^{\Pi},
  \tag{2.2.4.1}
\end{equation}

\VerdierSemanticAnchor{ch1-2-2-4-eq2}
\begin{equation}
  \int_{\operatorname{id}}^{\Pi}=\operatorname{id}.
  \tag{2.2.4.2}
\end{equation}

\VerdierSemanticAnchor{ch1-2-2-4-editorial-caveat}
编者在这里痛苦地不得不再次搬出那个惯常而恼人的把戏。公式 (2.2.4.1) 中的
等号只有在具有满足适当结合律性质的乘积函子时才成立，而实践中从来不是如此。

\VerdierSemanticAnchor{ch1-2-2-4-canonical-isomorphism}
一般而言，必须把“$=$”替换为满足 (1.2.1.3) 型通常相容性条件的典范同构。

\VerdierSemanticAnchor{ch1-2-2-4-sigma-analogue}
函子 $\int_{\varphi}^{\Sigma}$ 满足与 (2.2.4.1) 类似的关系。

\VerdierSemanticAnchor{ch1-2-2-4-summable-restriction}
当可数乘积与直和未必可表示时，把范围限制在 $\psi\varphi$-可求和复形上，
公式 (2.2.4.1) 仍成立。

\VerdierSemanticAnchor{ch1-2-2-4-simple-complex}
当 $\psi$ 是唯一的态射 $[n]\to[1]$ 时，采用记号 $\int_n$。对每个可求和的
$n$-重复形 $Y$，复形 $\int_nY$ 称为\emph{与 $Y$ 关联的单复形}。
